% ============================================================
%  CAPÍTULO 1: FUNDAMENTOS TEÓRICOS
% ============================================================
\chapter{Fundamentos teóricos}
\label{cap:fundamentos}

En este capítulo se presentan los fundamentos físicos y teóricos
sobre los que se sustenta el sistema fotoacústico desarrollado
en este trabajo. Se describe el efecto fotoacústico, su manifestación
en líquidos de baja absorción óptica, la necesidad de caracterizar
este tipo de medios y las técnicas actuales empleadas con ese fin.

% ============================================================
\section{El efecto fotoacústico}
\label{sec:efecto_fotoacustico}
% ============================================================

\subsection*{Origen histórico}

El efecto fotoacústico fue descrito por primera vez por Alexander
Graham Bell en 1880 \cite{Bell1880}. Bell trabajaba entonces, junto
con Charles Sumner Tainter, en el desarrollo del \textit{fotófono},
un instrumento con el que pretendía transmitir la voz a distancia
empleando la luz solar como portadora de la información
\cite{Marin2008}. Durante esos experimentos observó que una muestra
iluminada con radiación luminosa interrumpida periódicamente, y
colocada dentro de una celda cerrada, emitía un sonido audible.

El fenómeno permaneció prácticamente sin aplicaciones durante casi
un siglo, en buena medida porque no existían fuentes luminosas ni
detectores con la sensibilidad necesaria. Su reactivación ocurre en
la década de 1970 con los trabajos de Harshbarger y Robin, y de
Rosencwaig, que reintroducen la técnica como herramienta
espectroscópica; poco después, Rosencwaig y Gersho publican el
modelo analítico que describe la dependencia de la señal con la
frecuencia de modulación, el coeficiente de absorción y las
propiedades térmicas de la muestra \cite{RosencwaigGersho1976},
modelo que se convierte en la referencia estándar de la técnica
convencional \cite{PatelTam1981}.

\subsection*{Descripción del mecanismo físico}

El efecto fotoacústico consiste en la generación de una onda de
presión en un medio como consecuencia de la absorción de radiación
electromagnética. La cadena de conversión de energía comprende tres
etapas sucesivas \cite{Marin2008}:

\begin{enumerate}
    \item \textbf{Absorción óptica.} La muestra absorbe una fracción
    de la radiación incidente, determinada por su coeficiente de
    absorción óptica $\alpha$ a la longitud de onda de excitación.

    \item \textbf{Relajación no radiativa.} La energía absorbida
    promueve al medio a estados excitados que decaen predominantemente
    por vías no radiativas, depositando la energía como calor local.
    Esta conversión ocurre en un tiempo característico
    $\tau_\text{NR}$ mucho menor que las escalas temporales de
    interés en el experimento \cite{PatelTam1981}.

    \item \textbf{Expansión termoelástica.} El calentamiento local
    produce una dilatación del volumen iluminado; dicha dilatación
    actúa como fuente acústica y emite una onda de presión que se
    propaga por el medio hasta el transductor.
\end{enumerate}

La magnitud de la señal depende, por tanto, tanto de propiedades
ópticas ($\alpha$) como termodinámicas y mecánicas del medio
(coeficiente de expansión volumétrica $\beta$, calor específico a
presión constante $C_p$, densidad $\rho$ y velocidad del sonido
$v_s$). Esta doble dependencia es precisamente lo que convierte al
efecto fotoacústico en una técnica de caracterización: la señal
acústica contiene información que no es accesible por medición
óptica directa.

\subsection*{Régimen modulado y régimen pulsado}

Históricamente la técnica se desarrolló en el régimen de
\textit{excitación modulada}, en el que un haz continuo se interrumpe
periódicamente a frecuencias de audio y un micrófono detecta el
calentamiento y enfriamiento de una capa de gas en contacto térmico
con la muestra. Patel y Tam señalan que esta configuración es en
rigor una técnica \textit{fototérmica} más que fotoacústica, ya que
la onda acústica generada en la propia muestra desempeña un papel
menor y la detección depende de la difusión térmica hacia el gas,
un proceso poco eficiente; en consecuencia, su sensibilidad resulta
limitada y su aplicación práctica queda restringida a absorciones
relativamente altas \cite{PatelTam1981}.

El \textit{régimen pulsado}, empleado en este trabajo, sustituye la
modulación por pulsos láser cortos y el micrófono de gas por un
transductor piezoeléctrico en contacto directo con la muestra
líquida. La detección es entonces auténticamente fotoacústica: el
pulso de presión generado en el medio se registra directamente, sin
intermediarios térmicos, y el buen acoplamiento de impedancia
acústica entre líquido y transductor eleva la sensibilidad varios
órdenes de magnitud \cite{PatelTam1981}.

Independientemente del régimen considerado, la amplitud de presión
generada admite la forma
\begin{equation}
    p_0 = \text{const} \cdot \frac{\beta\, v_s^2}{C_p}\, E_0\, \alpha,
    \label{eq:presion_general}
\end{equation}
donde $p_0$ es la amplitud de presión acústica, $E_0$ la energía del
pulso láser, $\beta$ el coeficiente de expansión térmica, $C_p$ el
calor específico a presión constante, $v_s$ la velocidad acústica y
$\alpha$ el coeficiente de absorción óptica del medio
\cite{PatelTam1981}. El agrupamiento $\beta v_s^2 / C_p$ se conoce
como parámetro de Grüneisen y sintetiza la eficiencia con la que un
medio convierte calor depositado en presión acústica.

De la \cref{eq:presion_general} se desprende la propiedad que
sustenta el uso cuantitativo de la técnica: normalizando la señal
detectada por la energía del pulso, la cantidad resultante es
directamente proporcional al coeficiente de absorción óptica del
medio \cite{PatelTam1981}. Esta es la razón por la que el arreglo
experimental descrito en el \cref{cap:arreglo} incorpora un
fotodiodo dedicado a monitorear la energía de cada disparo.

% ============================================================
\section{El efecto fotoacústico en líquidos de baja absorción}
\label{sec:fotoacustico_liquidos}
% ============================================================

\subsection*{Geometría de la fuente acústica}

El caso de los líquidos de baja absorción óptica constituye un
régimen distinto y requiere un tratamiento propio. Si el producto
$\alpha l$ entre el coeficiente de absorción y la longitud del
camino óptico satisface
\begin{equation}
    \alpha l \ll 1,
    \label{eq:condicion_debil}
\end{equation}
la energía absorbida se distribuye de manera aproximadamente
uniforme a lo largo de toda la trayectoria del haz dentro de la
muestra. La fuente acústica deja entonces de ser puntual y se
convierte en un cilindro delgado coincidente con el haz, que emite
una \textit{onda cilíndrica} \cite{PatelTam1981}.

Esta situación contrasta con la de los medios fuertemente
absorbentes, donde la energía se deposita en una capa superficial
delgada y la emisión es esencialmente esférica desde el punto de
incidencia \cite{PatelTam1981}. La distinción no es un detalle
geométrico menor: determina la forma temporal del pulso registrado,
su ley de decaimiento con la distancia y, en consecuencia, la
manera correcta de interpretar la señal capturada por el
transductor.

Bajo la aproximación adiabática --- es decir, despreciando la
difusión de calor fuera del volumen iluminado durante la duración
del pulso, hipótesis justificada porque la longitud de difusión
térmica en escalas de nanosegundos es varios órdenes de magnitud
menor que el radio del haz --- el desarrollo riguroso muestra que
el pulso de presión resultante consta de un lóbulo positivo seguido
de uno negativo, separados por un intervalo aproximadamente igual a
la duración del pulso láser \cite{PatelTam1981}. Esta firma bipolar
es el rasgo distintivo de la generación fotoacústica en volumen y
es la que se busca identificar en las señales adquiridas.

\subsection*{Sensibilidad alcanzable}

La motivación central del régimen pulsado con detección
piezoeléctrica es la medición de coeficientes de absorción muy
pequeños. Patel y Tam reportan capacidad demostrada para medir
coeficientes del orden de \SIrange{e-4}{e-5}{\per\centi\meter},
un intervalo inaccesible para la espectrofotometría de transmisión
convencional, y aplican la técnica a la determinación de los
espectros de absorción visible del agua y del agua pesada, así como
a los armónicos vibracionales altos de líquidos orgánicos
transparentes como el benceno \cite{PatelTam1981}.

Este es el argumento que justifica el uso de la técnica
fotoacústica sobre alternativas ópticas directas cuando el medio de
interés es esencialmente transparente a la longitud de onda de
excitación, como ocurre con el agua en la región visible.

\subsection*{Detección con hidrófono de banda ancha}

La configuración empleada en este trabajo tiene un antecedente
directo en el trabajo de Quan \textit{et al.}, quienes generaron
ondas acústicas en agua destilada y metanol mediante pulsos cortos
de un láser Nd:YAG a \SI{1064}{\nano\meter} y las detectaron con un
hidrófono piezoeléctrico de banda ancha (\SIrange{0.1}{20}{\mega\hertz}),
obteniendo buena concordancia entre predicción teórica y medición
\cite{Quan1994}.

Dicho trabajo señala además tres factores que condicionan la forma
de onda registrada y que resultan pertinentes para el diseño del
arreglo experimental: el tiempo de subida del detector queda
determinado por el espesor del transductor y la velocidad del
sonido en el material piezoeléctrico; un detector con dimensiones
laterales mayores que la longitud de onda del pulso degrada tanto
la resolución espacial como la temporal; y las dimensiones finitas
de la celda que contiene el líquido dan lugar a reflexiones
acústicas múltiples que se superponen a la señal de interés
\cite{Quan1994}. Este último punto es determinante para el diseño
de la celda fotoacústica descrita en el \cref{cap:arreglo}.

\subsection*{Mecanismos competidores: electrostricción}

Conforme el coeficiente de absorción disminuye, la contribución
termoelástica descrita en la \cref{sec:efecto_fotoacustico} deja de
ser el único mecanismo capaz de generar presión en el medio. En
dieléctricos de baja absorción, la interacción entre el campo
electromagnético y los dipolos inducidos comprime el material hacia
las regiones de mayor intensidad de campo; este fenómeno se conoce
como \textit{electrostricción} y constituye un efecto inherentemente
transitorio, cuya escala temporal de compensación corresponde al
tiempo que tarda la onda acústica generada en recorrer un ancho de
haz \cite{Astrath2023}.

Astrath \textit{et al.} demuestran experimentalmente en agua que,
bajo excitación con haces pulsados fuertemente enfocados, la
electrostricción puede ser el efecto transitorio dominante en la
generación de la onda acústica. Su evidencia más contundente
proviene de mediciones realizadas a \SI{4.0}{\celsius}: a esa
temperatura el coeficiente de expansión térmica del agua se anula,
la contribución termoelástica se reduce drásticamente y la señal
remanente puede atribuirse inequívocamente a la electrostricción
\cite{Astrath2023}.

Este resultado tiene una consecuencia metodológica directa sobre el
presente trabajo. La dependencia de la señal fotoacústica con la
temperatura de la muestra no proviene únicamente de la variación de
$\beta$, $C_p$ y $v_s$ con $T$; también refleja el peso relativo de
mecanismos de generación distintos. Por ello, la capacidad de
registrar simultáneamente la señal acústica y la temperatura de la
muestra --- función central del sistema de automatización
desarrollado --- no es un accesorio del experimento, sino una
condición para interpretar correctamente los datos obtenidos.

% ============================================================
\section{Necesidad de caracterizar líquidos}
\label{sec:necesidad_caracterizacion}
% ============================================================

La caracterización de líquidos consiste en la determinación
cuantitativa de las propiedades que describen su comportamiento
físico: coeficiente de absorción óptica, difusividad térmica,
velocidad del sonido, coeficiente de expansión volumétrica y su
dependencia con la temperatura, entre otras.

Las técnicas fotoacústicas pulsadas se han aplicado a la
determinación de trazas de analitos en sistemas gaseosos y
líquidos, al ensayo no destructivo mediante ultrasonido inducido
ópticamente y al estudio de tiempos de vida de estados excitados
\cite{Quan1994}. En el ámbito de la espectroscopía de sistemas
acuosos, la configuración habitual consiste en disolver el analito
de interés en un medio de baja absorción de fondo, de modo que la
señal fotoacústica quede dominada por la especie que se desea
cuantificar \cite{Quan1994}.

La técnica ha encontrado también aplicación en la caracterización
de materiales y en problemáticas ambientales, campo en el que su
sensibilidad a absorciones pequeñas resulta particularmente
ventajosa \cite{Marin2008,Tam1986}.

Un aspecto relevante para este trabajo es que la señal fotoacústica
no depende exclusivamente de la absorción óptica. Como muestra la
\cref{eq:presion_general}, la amplitud registrada involucra también
$\beta$, $C_p$ y $v_s$, parámetros que son característicos de cada
líquido y que varían con la temperatura. En consecuencia, el
análisis conjunto de la amplitud, la forma temporal y el contenido
espectral de la señal permite extraer parámetros físicos del medio
que sirven para discriminar entre un conjunto acotado de líquidos
candidatos conocidos.

Conviene precisar el alcance de esta afirmación. La identificación
molecular abierta de una sustancia desconocida requiere realizar un
barrido espectral, esto es, excitar la muestra a múltiples
longitudes de onda y reconstruir su espectro de absorción, tal como
hicieron Patel y Tam para obtener los espectros del agua y del agua
pesada \cite{PatelTam1981}. Un láser de longitudes de onda fijas
como el empleado en este trabajo no permite ese barrido. El objetivo
perseguido aquí es, por tanto, la caracterización de parámetros
físicos y la discriminación entre candidatos conocidos, no la
identificación espectroscópica abierta.

% ============================================================
\section{Técnicas actuales de caracterización de líquidos}
\label{sec:tecnicas_actuales}
% ============================================================

\subsection*{Espectrofotometría de transmisión}

El método de referencia para determinar el coeficiente de absorción
óptica de un líquido es la medición directa de la transmitancia,
mediante espectrofotometría UV-Vis o infrarroja, y la aplicación de
la ley de Lambert-Beer, cuyos orígenes se remontan a los trabajos
de Lambert y Beer en los siglos XVIII y XIX \cite{Marin2008}.

Su limitación aparece precisamente en el régimen de interés de este
trabajo. Cuando $\alpha l \ll 1$, la diferencia entre la intensidad
incidente y la transmitida es del mismo orden que el ruido y las
pérdidas por reflexión e interfaces, de modo que la medición pierde
significado. Es esta limitación la que motivó el desarrollo de las
técnicas fotoacústicas para la medición de absorciones débiles
\cite{PatelTam1981}.

\subsection*{Fotoacústica de gas con micrófono}

La configuración convencional, derivada del trabajo de Bell y
reactivada en los años setenta, emplea un haz modulado y un
micrófono que detecta las variaciones de presión en una capa de gas
en contacto térmico con la muestra
\cite{RosencwaigGersho1976,PatelTam1981}. Su principal desventaja
para muestras condensadas es la dependencia de la difusión térmica
hacia el gas, un mecanismo de acoplamiento ineficiente que limita
la sensibilidad y restringe su utilidad a absorciones
relativamente altas \cite{PatelTam1981}.

\subsection*{Detección piezoeléctrica con excitación modulada}

Una mejora intermedia consiste en conservar la modulación a
frecuencias de audio pero sustituir el micrófono por un transductor
piezoeléctrico en contacto directo con la muestra sólida o líquida,
aprovechando el mejor acoplamiento de impedancia acústica
\cite{PatelTam1981}. La sensibilidad mejora respecto a la
detección con micrófono, pero la excitación modulada continúa
limitando la relación señal a ruido alcanzable frente al régimen
pulsado.

\subsection*{Métodos ópticos de detección}

La onda de presión también puede detectarse por vía óptica, sin
transductor en contacto. Las variaciones locales de presión alteran
el índice de refracción del medio, efecto que puede sondearse
mediante interferometría o técnicas de lente fototérmica
\cite{Astrath2023}; alternativamente, puede medirse la deflexión de
un haz de prueba que atraviesa la región perturbada \cite{Quan1994}.
Estos métodos evitan perturbar mecánicamente la muestra, a costa de
un arreglo óptico considerablemente más complejo y de una
alineación más delicada.

\subsection*{Fotoacústica pulsada con detección piezoeléctrica}

La técnica adoptada en este trabajo combina excitación por pulsos
láser cortos, detección piezoeléctrica en contacto directo con el
líquido y adquisición sincronizada con el disparo del láser
\cite{PatelTam1981}. Sus ventajas frente a las alternativas
anteriores son tres: sensibilidad a coeficientes de absorción
varios órdenes de magnitud menores que los accesibles por
transmisión; acceso a la forma temporal completa del pulso
acústico, y no solo a su amplitud, lo que habilita el análisis
espectral de la señal; y un arreglo experimental mecánicamente más
simple que el requerido por los métodos de detección óptica.

Su principal desventaja es de naturaleza operativa: cada punto
experimental exige coordinar el disparo del láser, la adquisición
del osciloscopio, el monitoreo de la energía del pulso y el
registro de la temperatura de la muestra. Realizada manualmente,
esta coordinación es lenta, propensa a error e imposible de
reproducir con exactitud entre sesiones. Es esta limitación la que
el sistema de automatización descrito en los capítulos siguientes
se propone resolver.
